\chapter{The Chatman Equation}
\label{chap:chatman_equation}

\section{The Need for a Formal Statement}

The enterprise process gap (Chapter~\ref{chap:enterprise_gap}) identifies the problem. The 2016 experiment (Chapter~\ref{chap:origin}) identifies the failure mode. What was missing, for the first several years of this research program, was a formal statement of what a consequential action requires --- a statement precise enough to serve as an admissibility criterion, rigorous enough to support algebraic manipulation, and operational enough to be implemented in software.

The Chatman Equation is that formal statement. It is not a metaphor. It is not a slogan. It is an equation. It has a left side, a right side, and a receipted form. It can be evaluated. It can be falsified. It has been implemented across the thirty-four projects documented in this corpus, and it has been used to distinguish \textsc{alive} verdicts from \textsc{partial} ones with the precision that customer accountability requires.

\section{The Equation}

The Chatman Equation is:

\[
A = \mu(O^{*})
\]

Where:

\begin{itemize}
  \item $A$ is a consequential action --- an act that produces a durable, verifiable state change in a process.
  \item $O^{*}$ is the observed operational state: the complete set of process-event evidence available at the moment of action, including event logs, object states, agent assignments, and system boundary crossings.
  \item $\mu$ is the projection function: the transformation that converts observed operational state into an admissible basis for action.
\end{itemize}

The equation states that a consequential action is only valid when it is grounded in --- projected from --- observed operational state. An action that is not grounded in observed operational state is not $A$. It may be an assertion. It may be a prediction. It may be a hallucination. It is not a consequential action in the sense the equation defines.

\section{Unpacking the Components}

\subsection{Observation ($O^{*}$)}

$O^{*}$ is not a snapshot. It is a process-event history. In the formal language of process mining, it is an object-centric event log: a structured record of events, each with a timestamp, an activity label, a set of objects, and a set of attribute values. The OCEL 2.0 standard \cite{vanderaalst2022ocel} provides the canonical schema for $O^{*}$.

The asterisk in $O^{*}$ is significant. It denotes the complete observed history up to the moment of action --- not a sample, not a summary, not a model-generated approximation. If the history is incomplete, $O^{*}$ is incomplete, and the equation's right side is undefined. This is the formal basis for the \textsc{partial} verdict: when $O^{*}$ cannot be fully constructed from available evidence, the equation cannot be evaluated, and the action cannot be certified as consequential.

\subsection{Projection ($\mu$)}

The projection function $\mu$ is the transformation from observed state to admissible action basis. It is not a prediction. Prediction manufactures a plausible next token. Projection derives, from verified evidence, a valid basis for a specific action.

$\mu$ must be:
\begin{itemize}
  \item \textbf{Deterministic}: given the same $O^{*}$, $\mu(O^{*})$ must produce the same admissibility decision.
  \item \textbf{Auditable}: the transformation must be inspectable, not opaque.
  \item \textbf{Bounded}: $\mu$ must operate on a well-defined evidence scope, not unbounded context.
\end{itemize}

In the implementation, $\mu$ is realized through SPARQL queries against RDF-encoded process graphs (Chapter~\ref{chap:ggen}), through conformance-checking algorithms in \texttt{sources-wasm4pm}, and through the governance rules encoded in \texttt{blue-river-dam}. Each realization is a specific, inspectable instantiation of $\mu$ for a particular process domain.

\subsection{Action ($A$)}

$A$ is a consequential action. It has three required properties.

First, $A$ must produce a durable state change. A state change is durable if it is recorded in the event log and survives the termination of the process instance that produced it.

Second, $A$ must be attributable to a specific agent. Coordination requires agency. An action without an attributable agent is a process event without a responsible party --- a compliance gap.

Third, $A$ must be replayable. The event produced by $A$ must be capable of being replayed against the declared process model, producing a conformance score. If the event is not replayable, $A$ cannot be verified as lawful.

\section{The Receipted Form}

The base equation $A = \mu(O^{*})$ establishes the admissibility condition. The receipted form extends it to include proof of satisfaction:

\[
R \vdash A = \mu(O^{*})
\]

Where $R$ is the receipt: a cryptographically bound record that certifies the equation was evaluated, the evaluation succeeded, and the action $A$ was admitted on the basis of projection from observed operational state $O^{*}$.

The receipt $R$ is not a hash of the output. It is a proof of the evaluation. It records:
\begin{itemize}
  \item The identity of the observation $O^{*}$ (by content-addressable reference).
  \item The identity of the projection function $\mu$ applied.
  \item The identity of the action $A$ admitted.
  \item The timestamp of evaluation.
  \item A cryptographic commitment linking all of the above.
\end{itemize}

The receipted form is the basis for the \textsc{alive} verdict. A project is \textsc{alive} when, for each claimed consequential action in its scope, $R \vdash A = \mu(O^{*})$ can be evaluated and the receipt can be verified. A project is \textsc{partial} when some actions are receipted and others are not. A project is \textsc{blocked} when the receipt chain cannot be constructed --- when $O^{*}$ is unavailable, $\mu$ is undefined, or $A$ cannot be verified.

\section{Admissibility and Consequence}

The equation introduces the concept of admissibility: an action is admissible when its basis in observed operational state can be demonstrated. Admissibility is the process-intelligence analog of legal admissibility in evidence law. A claim that cannot be traced to observable evidence is not inadmissible because it is false. It is inadmissible because it cannot be verified. The distinction matters: inadmissibility is a structural property of the claim, not a judgment about its content.

This framing has significant consequences for how the equation is applied in practice. When an AI system asserts that a process has completed --- that a handoff has occurred, that an artifact is in a particular state, that a compliance control has been applied --- the assertion is not evaluated by asking whether the assertion is plausible. It is evaluated by asking whether $R \vdash A = \mu(O^{*})$ can be verified for the claimed action. If the receipt does not exist, the claim is \textsc{partial}, regardless of how plausible the assertion sounds.

This is the formal basis for the research program's refusal of hallucination as an acceptable failure mode. Hallucination is not a technical artifact of model behavior. It is a structural consequence of applying prediction systems in coordination roles, where admissibility requirements exist and cannot be satisfied by prediction alone.

\section{Relation to the Research Corpus}

The Chatman Equation is not a theoretical construct that was developed independently of the research projects. It was derived from, and verified against, the projects. Each ALIVE verdict in the checkpoint ledger (\texttt{checkpoints/}) is a verified instance of $R \vdash A = \mu(O^{*})$. Each PARTIAL verdict is a documented case where the equation's components are partially satisfied. The thirty-four projects are, collectively, the empirical basis from which the equation was derived and against which it has been evaluated.

The equation also provides the formal spine of Chapters~\ref{chap:process_evidence} through~\ref{chap:conclusion}. Every subsequent chapter can be read as a detailed account of one or more components of the equation: how $O^{*}$ is constructed (Chapter~\ref{chap:process_evidence}), how $\mu$ is implemented (Chapter~\ref{chap:ggen}), how $A$ is constrained by grammar (Chapter~\ref{chap:command_grammar}), and how $R$ grounds capital-market claims (Chapter~\ref{chap:ai_xynz}).
